\chapter{FastTree}
\section{Introduction}
Existing LLM serving systems use a radix tree to organize the global key-value (KV) cache, facilitating cache reuse across different queries and thus reducing unnecessary memory use. Despite this, these systems still rely on conventional computation patterns for attention operations, resulting in redundant memory loads and GPU tensor core underutilization.

FastTree introduces GPU kernels tailored for efficiently processing queries that share contexts through the radix tree.

\begin{figure}[H]
    \centering
    \includegraphics[width=0.5\textwidth]{\FIGDIR/kv_cache_sharing_through_radix_tree.png}
    \caption{KV cache sharing through Radix Tree }
\end{figure}

\begin{figure}[H]
    \centering
    \includegraphics[width=0.5\textwidth]{\FIGDIR/computation_pattern_existing_system.png}
    \caption{Computation pattern of existing systems}
\end{figure}


However, although optimizing the memory layouts with a radix tree, existing systems still utilize the conventional attention kernels which concatenate the KV cache for each query and dispatch its computation into an individual GPU thread block. It requires redundantly loading the shared KV cache from slow global memory for every query and cannot meet the minimum shape requirements for tensor cores without padding.


With a given radix tree of the KV cache
at runtime, we can generate enormous grouping plans by changing the shared contexts used for query aggregation, as the different queries can share varied prefixes. Moreover, the plan selection can significantly affect the execution performance of the attention kernel, as different plans can introduce distinct overheads, including wasted shared memory and computation due to tile padding and IO cost due to increased intermediate results.

